\documentclass[a4paper]{amsart}
\usepackage{amssymb, mathtools, amsthm}
% \usepackage{a4wide}
\usepackage[hidelinks]{hyperref}
\usepackage{tikz}
\usepackage{graphicx}
\usepackage{todonotes}
\usepackage[font=small]{caption}
\usetikzlibrary{decorations.markings}

\newcommand{\authornote}[1]{\ignorespaces\marginpar{
\tiny #1}\ignorespaces}
\newcommand{\CP}{\mathbf{CP}}
\newcommand{\cO}{\mathcal{O}}
\newcommand{\FS}{\text{FS}}
\newcommand{\uS}{\mathbb{S}}
\newcommand{\sph}{{\uS^2}}
\newcommand{\laplace}{\Delta}
\newcommand{\Hess}{\nabla^2}
\newcommand{\Var}{\operatorname{Var}}
\newcommand{\dist}{\operatorname{dist}}
\newcommand{\trace}{\operatorname{tr}}
\newcommand{\re}{\operatorname{Re}}
\newcommand{\R}{\mathbb{R}}
\newcommand{\E}{\mathbb{E}}
\newcommand{\ii}{{\rm i}}
\newcommand{\dd}{{\rm d}}
\newcommand{\pv}{\operatorname{P.V.}}

%%%%%%%%% Theorem-like environment %%%%%%%%%%%
\theoremstyle{plain} %text of this environment is typesetted in italics
\newtheorem*{theorem}{Theorem}
\newtheorem*{lemma}{Lemma}
% \newtheorem{corollary}[theorem]{Corollary}
% \newtheorem{proposition}[theorem]{Proposition}
\newtheorem*{conjecture}{Conjecture}
\newtheorem*{problem}{Problem}
%
\theoremstyle{definition} %text of this environment is typesetted in roman letters
\newtheorem{definition}{Definition}
\newtheorem{assumption}{Assumption}

\theoremstyle{remark}
\newtheorem{remark}{Remark}
\newtheorem{example}{Example}

\begin{document}

\title[D and P locally minimize Gauss variance]{The Primitive and Diamond surfaces locally minimize\\ the variance of Gauss curvature}

\author{Hao Chen}
\email{chenhao5@shanghaitech.edu.cn}
\address{ShanghaiTech University, Shanghai, China}

\keywords{minimal surfaces}
\subjclass[2010]{Primary 53A10}

\date{\today}

\begin{abstract}
	We prove that Schwarz' Primitive and Diamond surfaces are local minimizers
	for the variance of Gauss curvature among local deformations within the space
	of Triply Periodic Minimal Surfaces of genus 3.  This is done by viewing the
	branched values of the Gauss map as a point configuration on the sphere, and
	rewriting the variance of Gauss curvature as a product of two integrals
	involving exponentials of Green functions.
\end{abstract}

\maketitle

Triply Periodic Minimal Surfaces (TPMSs) are minimal surfaces in flat 3-tori.
The genus of a TPMS is at least three.  A TPMS of genus three (TPMSg3) has a
Weierstrass parameterization defined on a branched double cover of the sphere
with eight branch points.  The branched covering map is provided by the Gauss
map, so the branch points correspond to eight \emph{flat points} on the TPMS
(where the Gauss curvature vanishes).  We refer the readers to~\cite{meeks1990}
for general reference.

TPMSs find applications in natural sciences as a model for periodic
bicontinuous structures~\cite{hyde1996}.  In nature and in laboratories, the
chance to observe Schwarz' Primitive (P), Diamond (D) and Schoen's Gyroid (G)
surfaces is much higher than finding other TPMSs.  Note that these surfaces
belong to the same Bonnet family.  They are therefore locally isometric, have
the same branch points, and the same variance of Gauss curvature.  For all the
three surfaces, the branch points are vertices of a cube.  Based on numerical
evidence~\cite{schroderturk2006}, it was conjectured that

\begin{conjecture}
	Schwarz' Primitive, Diamond, and Schoen's Gyroid surfaces have the minimal
	variance of Gauss curvature among all Triply Periodic Minimal Surfaces of
	genus 3.
\end{conjecture}

In this short note, we prove that the Primitive and Diamond are indeed local
minimizers.

\begin{theorem} \label{thm:main}
	Schwarz' Primitive and Diamond surfaces are local minimizers for the variance
	of Gauss curvature among local deformations within the space of Triply
	Periodic Minimal Surfaces of genus 3.
\end{theorem}

We will relax the problem to study the variance for the Gauss curvature defined
on the $13$-dimensional space $V$ of configurations of eight points on the
sphere up to global rotations.  They correspond to the normal vectors at the
eight flat points on a TPMSg3.  For Diamond, Primitive, and Gyroid surfaces,
the eight points are at the vertices of a cube.  However, only a 5-dimensional
variety in $V$ correspond to TPMSg3s.  In particular, the 5-dimensional
subspace of antipodal configurations correspond to embedded TPMSg3s known as
the \emph{Meeks family}.

In~\cite{koiso2018}, it was proved that, up to homotheties, the Primitive,
Diamond, and Gyroid all belong to a unique locally rigid 5-parameter smooth
family of pairwise non-homothetic triply periodic minimal surfaces.  In
particular, the family that contains the Diamond and the Primitive surfaces is
the Meeks family.  Schoen's Gyroid, however, does not belong to Meeks family.

At the cubic configuration, we compute the Hessian of the variance for the
Gauss curvature on $V$ and find, unfortunately, a $2$-dimensional negative
eigenspace generated by ``twists''.  So the Hessian is indefinite, and the
cubic configuration is not a local minimizer in $V$.  Nevertheless, we are able
to prove positive definiteness for the deformations preserving antipodality, so
the cubic configuration is a local minimizer among antipodal deformations,
thereby proving the Theorem.

The $5$-parameter family of local deformations of the Gyroid is not fully
understood.  Among the known deformations of the Gyroid, it is indeed a local
minimizer.  But in the presence of the negative eigenspace, we can not rule out
a possibility that it might not be a local minimizer along other deformations.

\subsection*{Acknowledgement}

The author thanks Chengjian Yao and Thierry De Pauw for very helpful and
inspiring discussions.

\section{Reformulation}

\subsection{Weierstrass parameterization}

Given a TPMSg3, let $p_1, \cdots, p_8$ be (the stereographic projection of) the
Gauss map at the eight flat points.  Consider the hyperelliptic Riemann surface
$\Sigma$ of genus three defined by
\[
	w^2 = P(z) = \prod_{i=1}^8 (z-p_i).
\]
It can be seen as a branched double cover of the Riemann sphere $\uS^2 \simeq
\CP^1$ with eight branch points.  Then we have the following \emph{Weierstrass
parameterization} for the TPMS (up to translations)
\begin{equation}\label{eq:weierstrass}
	\Sigma \ni (z, w) \mapsto \re \int^{(z, w)} \frac{(1-z^2, \ii(1+z^2), 2z)}{w}
	e^{i \vartheta}\, dz.
\end{equation}
Switching the two branches by $(z,w) \mapsto (z,-w)$ correspond to the
inversion with respect to any of the eight flat points.  The parameter
$\vartheta$ is the \emph{Bonnet angle}; changing $\vartheta$ gives a locally
isometric deformation of the surface known as \emph{Bonnet rotation}.  Two
surfaces whose $\vartheta$ differ by $\pi/2$ are called \emph{conjugate pairs}.

For the minimal surface to be triply periodic and globally well-defined, one
needs to solve the period problem (aka, \emph{close the periods}).  More
precisely, one finds $p_i$'s and $\vartheta$ so that the real parts of the
integrals over closed cycles on $\Sigma$ form a three dimensional lattice.  The
period problem involves three real equations for each of the three coordinates,
hence nine real equations in total.  Since we have $17$ real parameters
($p_i$'s and $\vartheta$), the set of TPMSg3s form a $8$-dimensional variety up
to translations, or $5$-dimensional variety up to homotheties.

Closing the periods is usually very difficult.  But Meeks proved
in~\cite{meeks1990} that, when the branch points form four antipodal pairs,
there exists a conjugate pair embedded TPMSg3s.  Up to homotheties, these
TPMSg3s form a $5$-dimensional manifold termed \emph{Meeks family}.

In this paper, we will not insist on closing the periods.  The variance of
Gauss curvature is invariant under local isometries such as Bonnet rotations.
So it is a function defined on a $13$-dimensional space of configurations of
eight points on the sphere up to rotations.  The Meeks surfaces correspond to
the 5-dimensional subspace of antipodal configurations, for which we know the
existence of $\vartheta$ that closes the periods.  Our main Theorem is
equivalent to that the cubic configuration is a local minimizer of the Gauss
variance on this $5$-dimensional subspace.

\subsection{Variance of Gauss curvature}

More precisely, we want to minimize the (dimensionless) variance of Gauss
curvature, defined by
\[
	\sigma(K)
	= \frac{\langle K^2 \rangle - \langle K \rangle^2}{\langle K \rangle^2}
	= \Bigg[\frac{\int_\Sigma K^2 \, dA}{\int_\Sigma dA} - \Bigg( \frac{\int_\Sigma K \, dA}{\int_\Sigma dA} \Bigg)^2\Bigg] / \Bigg( \frac{\int_\Sigma K \, dA}{\int_\Sigma dA} \Bigg)^2.
\]
By Gauss-Bonnet, $\int_\Sigma K \, dA$ is a topological invariant.  It is then
equivalent to minimize
\begin{equation}\label{eq:form1}
	\int_\Sigma K^2 dA \times \int_\Sigma dA.
\end{equation}

From the Weierstrass parameterization, we compute
\[
	dA = \frac{(1+|z|^2)^2}{|w|^2} \, dx \wedge dy, \qquad K = -4 \frac{|w|^2}{(1+|z|^2)^4},
\]
where $z = x + y \ii$.  Indeed, one verifies that
\[
	\int_\Sigma K dA = -4 \int_\Sigma \frac{dx \wedge dy}{(1+|z|^2)^2} = - 8 \pi
\]
as expected.  The area element and the Gauss curvature are invariant after
switching the branches.  It then suffices to compute the invariance on one of
the branch.  So \eqref{eq:form1} becomes (up to a constant factor)
\begin{equation}\label{eq:form2}
	\int_{\CP^1} \frac{|P(z)|}{(1+|z|^2)^6} \, dx \wedge dy \times
	\int_{\CP^1} \frac{(1+|z|^2)^2}{|P(z)|} \, dx \wedge dy.
\end{equation}
Note that the second integral is improper: The integrand is not defined at the
zeros of $P(z)$.  So the integral is defined as the limit
\[
	\int_{\CP^1} \frac{(1+|z|^2)^2}{|P(z)|} \, dx \wedge dy
	= \lim_{\varepsilon \to 0}
	\int_{\CP^1 \setminus \cup_i D^i_\varepsilon} \frac{(1+|z|^2)^2}{|P(z)|} \, dx \wedge dy,
\]
where $D^i_\varepsilon$ denote the $\varepsilon$-disks around $p_i$ under the
Fubini--Study metric.  Fortunately, the limit converges.

\subsection{Fubini--Study metric}

Note that
\[
	\dd A_\FS = \frac{4 \, dx \wedge dy}{(1+|z|^2)^2}
\]
is the area element for the Fubini--Study metric on $\CP^1$, which is
(normalized to) exactly the standard metric on the unit sphere $\uS^2$.
Moreover, we can identify $P(z)$ to a global section $s$ of the holomorphic
line bundle $\cO(8) = \cO(1)^{\otimes 8}$ over $\CP^1$, and
\[
	\|s\|_\FS = \frac{|P(z)|}{(1+|z|^2)^4}
\]
is its point-wise Fubini--Study norm.  See, for instance,
\cite{huybrechts2005}.  Then \eqref{eq:form2} becomes (up to a constant factor)
\begin{equation}\label{eq:form3}
	J_8 := 
	\int_{\CP^1} \| s \|_\FS \, \dd A_\FS \times
	\int_{\CP^1} \| s \|_\FS^{-1} \, \dd A_\FS. 
\end{equation}
Note again that the second integral is improper.  It is defined as the
convergent limit of the integral over $\CP^1 \setminus \cup_i D^i_\varepsilon$
as $\varepsilon \to 0$.

\subsection{Lelong--Poincar\'e Equation}

Define
\[
	u = \log \| s \|_\FS.
\]
Note that the Fubini--Study metric on $\cO(8)$ is Hermitian.  One then easily verifies the
Lelong--Poincar\'e Equation
\[
	\laplace u = 2\pi \sum_{i=1}^8 \delta(z-p_i) - 4\frac{4}{(1+|z|^2)^2},
\]
or, using the Laplacian under the Fubini--Study metric,
\[
	\laplace_\FS u = \frac{(1+|z|^2)^2}{4} \laplace u
	= 2\pi \sum_{i=1}^8 \delta_i(z) - 4,
\]
where
\[
	\delta_i(z) = \frac{(1+|z|^2)^2}{4} \delta(z-p_i)
\]
denote the Dirac delta function under the Fubini--Study metric centered at
$p_i$.

The solution is given by
\[
	u = 2\pi \sum_{i = 1}^8 G_i(z),
\]
where $G_i(z)$ is the Green function solving
\begin{equation}\label{eq:laplaceG}
	\laplace_\FS G_i(z) = \delta_i(z) - \frac{1}{4\pi}.
\end{equation}
On the unit sphere $\uS^2 \simeq \CP^1$, we have
\[
	G_i(z) = -\frac{1}{2\pi} \log \sin \frac{\dist(z, p_i)}{2} + C
\]
where $\dist(z, p_i)$ is the spherical distance between $z$ and $p_i$.  Then
\eqref{eq:form3} becomes
\begin{equation}\label{eq:form4}
	J_8 =
	\int_{\uS^2} \exp(u) \, \dd A_\sph \times
	\int_{\uS^2} \exp(-u) \, \dd A_\sph,
\end{equation}
where $\dd A_\sph$ is the standard spherical area element on the unit sphere.
Note that \eqref{eq:form4} does not depend on the constant $C$ in $G_i$, so we
assume $C=0$ from now on.  Note again that the second integral is improper and
is defined as the convergent limit
\[
	\int_{\uS^2} \exp(-u) \, \dd A_\sph = \lim_{\varepsilon \to 0}
	\int_{\Omega_\varepsilon} \exp(-u) \, \dd A_\sph,
\]
where
\[
	\Omega_\varepsilon = \uS^2 \setminus \cup_{i=1}^n D^i_\varepsilon.
\]

The Conjecture would be proved if $J_8$ is minimized when $p_i \in \uS^2$ are
the vertices of a cube.

\begin{remark}
	With this reformulation, we actually expand our consideration beyond TPMSg3s.
	We are now considering all minimal surfaces with the Weierstrass
	parameterization~\eqref{eq:weierstrass}.  In particular, we do not need the
	period problems to be solved, not to mention the embeddedness.
\end{remark}

\subsection{Generalization}

The analysis above can be generalized as follows.  Fix $n$ points $p_1, \cdots, p_n$, and define the polynomial
\[
	P(z) = \prod_{i=1}^n (z-p_i).
\]
Consider the product of integrals
\[
	\int_{\CP^1} \frac{|P(z)|}{(1+|z|^2)^{\frac{n}{2}+2}} \, dx \wedge dy \times
	\int_{\CP^1} \frac{(1+|z|^2)^{\frac{n}{2}-2}}{|P(z)|} \, dx \wedge dy.
\]
Note that
\[
	\|s\|_\FS = \frac{|P(z)|}{(1+|z|^2)^{n/2}}
\]
is the point-wise Fubini--Study norm of $P(z)$ as a global section $s$ of
$\cO(n)$ over $\CP^1$. Then the product can be rewritten (up to a constant
factor) into the form
\[
	J_n :=
	\int_{\CP^1} \| s \|_\FS \, \dd A_\FS \times
	\int_{\CP^1} \| s \|_\FS^{-1} \, \dd A_\FS. 
\]
Now the Lelong--Poincar\'e Equation is
\[
	\laplace_\FS u = 2\pi \sum_{i=1}^n \delta_i(z) - n/2,
\]
which is solved by
\[
	u = 2\pi \sum_{i=1}^n G_i(z).
\]
So
\[
	J_n = 
	\int_{\uS^2} \exp(u) \, \dd A_\sph \times
	\int_{\uS^2} \exp(-u) \, \dd A_\sph.
\]
And we could ask the following problem
\begin{problem}
	What are the positions of $p_1, \cdots, p_n$ that minimize $J_n$?
\end{problem}
When $n = 8$, this is exactly the problem of finding branch values to minimize
the variance of the Gauss curvature.

\section{Gradient and Hessian}

\subsection{The gradient}

Recall that
\[
	\Omega_\varepsilon = \uS^2 \setminus \cup_{i=1}^n D^i_\varepsilon,
\]
where $D^i_\varepsilon$ denote the $\varepsilon$-disks around $p_i$ under the
spherical metric.  Define
\[
	I^\pm_\varepsilon = \int_{\Omega_\varepsilon} \exp(\pm u) \, \dd A_\sph,
\]
and two probability measure $\dd\mu^\pm_\varepsilon = \rho^\pm_\varepsilon \dd
A_\sph$ on $\Omega_\varepsilon$, where
\[
	\rho_\varepsilon^\pm
	= \frac{\exp(\pm u)}{\int_{\Omega_\varepsilon} \exp(\pm u) \, \dd A_\sph}
	= \frac{\exp(\pm u)}{I^\pm_\varepsilon}.
\]

We compute
\begin{align*}
	\langle \nabla_{\! p_i} \log I^\pm_\varepsilon, v_i\rangle &=
	\frac{\langle \nabla_{\! p_i} I^\pm_\varepsilon, v_i\rangle}{I^\pm_\varepsilon}
	\\ &=
	\frac{1}{I^\pm_\varepsilon}
	\Big(
		\int_{\Omega_\varepsilon} \langle \pm 2\pi \nabla_{\! p_i} G_i, v_i \rangle \exp(\pm u) \, \dd A_\sph
		+ \oint_{\partial D^i_\varepsilon} \exp(\pm u) \langle n_i, \tilde v_i \rangle \dd s
	\Big)
	\\ &=
	\int_{\Omega_\varepsilon} \langle \pm 2\pi \nabla_{\! p_i} G_i, v_i \rangle \rho_\varepsilon^\pm \, \dd A_\sph
	+ \oint_{\partial D^i_\varepsilon} \langle n_i, \tilde v_i \rangle \rho_\varepsilon^\pm \dd s
\end{align*}
where $v_i \in T_{p_i} \uS^2$, $n_i$ is the unit normal vectors on the boundary
of $D^i_\varepsilon$ pointing to the interiors of the disks, and $\tilde v_i$
is a Killing field such that $\tilde v_i (p_i) = v_i$. So for
\[
	v = (v_1, \cdots, v_8) \in \bigoplus_{i=1}^8 T_{p_i} \uS^2,
\]
we have
\begin{equation}\label{eq:grad}
	\langle \nabla_{\! p} \log I^\pm_\varepsilon, v\rangle 
	= \int_{\Omega_\varepsilon} \sum_{i=1}^n \langle \mp 2\pi \nabla_{\! z} G_i, \tilde v_i \rangle \dd\mu_\varepsilon^\pm
	+ \sum_{i=1}^n \oint_{\partial D^i_\varepsilon} \langle n_i, \tilde v_i \rangle \rho_\varepsilon^\pm \dd s
\end{equation}

Note that the formula does not depend on the choice of the Killing fields
$\tilde v_i$.  To see this, assume another choice $\tilde v'_i$.  Then $\tilde
w_i = \tilde v'_i - \tilde v_i$ is a Killing field that vanishes at $p_i$,
which is given as a rotation around $p_i$.  Consequently, we have $\langle \mp
2\pi \nabla_{\! z} G_i, \tilde w_i \rangle = 0$ and $\langle n_i, \tilde w_i
\rangle = 0$.

Note also that, when $\tilde v_i = \tilde v_*$ is the same Killing field for all
$i$, then
\begin{align*}
	\langle \nabla_{\! p} \log I^\pm_\varepsilon, v\rangle 
	&= \int_{\Omega_\varepsilon} \sum_{i=1}^n \langle \mp 2\pi \nabla_{\! z} G_i, \tilde v_* \rangle \dd\mu_\varepsilon^\pm
	+ \sum_{i=1}^n \oint_{\partial D^i_\varepsilon} \langle n_i, \tilde v_* \rangle \rho_\varepsilon^\pm \dd s\\
	&= \int_{\Omega_\varepsilon} \langle \mp \nabla_{\! z} u, \tilde v_* \rangle \dd\mu_\varepsilon^\pm
	+ \int_{\Omega_\varepsilon} \langle \pm \nabla_{\! z} u, \tilde v_* \rangle \dd \mu_\varepsilon^\pm = 0
\end{align*}
So a deformation along a Killing field does not change $I^\pm_\varepsilon$.

\medskip

Recall that $2\pi G_i = \log \sin (\frac{1}{2} \dist(z, p_i))$, so
\[
	2\pi \nabla_{\! z} G_i \sim \frac{\hat w_i}{\dist(z, p_i)}, \qquad
	\rho^\pm_\varepsilon \sim \dist(z, p_i)^{\pm 1} \qquad
\]
as $z \to p_i$, where $\hat w_i = \nabla_z \dist(z, p_i)$.  It is then
obvious that~\eqref{eq:grad} converges as $\varepsilon \to 0$ for
$\mu^+_\varepsilon$.  As for $\mu^-_\varepsilon$, in the standard spherical
coordinate with the pole $p_i$, polar angle $\theta$ and azimuth angle $\phi$,
the integrands
\begin{align*}
	\langle 2\pi \nabla_{\! z} G_i, \tilde v_i \rangle \dd\mu_\varepsilon^-
	&\sim \frac{2}{\theta I^-} \langle \hat w_i(\theta, \phi), \tilde v_i(\theta, \phi) \rangle \dd \theta \dd \phi,\\
	\langle n_i, \tilde v_i \rangle \rho_\varepsilon^- \dd s
	&\sim \frac{-2}{I^-} \langle \hat w_i(\varepsilon, \phi), \tilde v_i(\varepsilon, \phi) \rangle \dd \phi
\end{align*}
as $\theta, \varepsilon \to 0$.  The integrals in~\eqref{eq:grad} then converge
as $\varepsilon \to 0$ because
\[
	\lim_{\theta \to 0} \int_0^{2\pi} \langle \hat w_i(\theta, \phi), \tilde v_i(\theta, \phi) \rangle \dd \phi \to 0.
\]
More specifically, we have
\[
	\langle \nabla_{\! p} \log I^\pm, v\rangle :=
	\lim_{\varepsilon\to 0}\langle \nabla_{\! p} \log I^\pm_\varepsilon, v\rangle 
	= \int_{\uS^2} \sum_{i=1}^n \langle \mp 2\pi \nabla_{\! z} G_i, \tilde v_i \rangle \dd\mu^\pm.
\]
As a sanity check, one easily verifies that this vanishes when $\tilde v_i =
\tilde v_*$ is the same Killing field for all $i$.

\medskip

Finally, by symmetry, one easily verifies that $\langle \nabla_{\! p} \log
I^\pm_\varepsilon, v\rangle = 0$ for any $v$ when $p_1,\cdots,p_8$ are the
vertices of the cube. So the cube is indeed a critical point of $J_8$.

\subsection{The Hessian}

We compute,
\begin{align*}
	\langle \nabla_{\! p_i} \rho^\pm_\varepsilon, v_i\rangle
	&=
	\frac{\langle \pm 2\pi \nabla_{\! p_i} G_i, v_i \rangle \exp(\pm u)}{I_\varepsilon^\pm}
	-
	\frac{\exp(\pm u) \langle \nabla_{\! p_i} I_\varepsilon^\pm, v_i\rangle}{(I_\varepsilon^\pm)^2}
	\\
	&= \rho_\varepsilon^\pm \Big(
		\langle \mp 2\pi \nabla_{\! z} G_i, \tilde v_i \rangle
		- \int_{\Omega_\varepsilon} \langle \mp 2\pi \nabla_{\! z} G_i, \tilde v_i \rangle \dd\mu_\varepsilon^\pm
		- \oint_{\partial D^i_\varepsilon} \langle n_i, \tilde v_i \rangle \rho_\varepsilon^\pm \dd s
	\Big)
\end{align*}

The Hessian
\begin{align}
	\Hess_{\! p} \log I^\pm_\varepsilon (v , v)
	=&
	\langle \nabla_{\! p} \langle \nabla_{\! p} \log I^\pm_\varepsilon, v\rangle, v \rangle \nonumber \\
	=&\phantom{+}
	\int_{\Omega_\varepsilon} \sum_{i=1}^n \langle \nabla_{\! p_i} \langle \pm 2\pi \nabla_{\! p_i} G_i, v_i \rangle, v_i \rangle \rho_\varepsilon^\pm \, \dd A_\sph \label{eq:l1}\\
	&+ \int_{\Omega_\varepsilon} \sum_{i=1}^n \langle \pm 2\pi \nabla_{\! p_i} G_i, v_i \rangle \sum_{j=1}^n \langle \nabla_{\! p_j} \rho_\varepsilon^\pm, v_j \rangle \, \dd A_\sph \label{eq:l2}\\
	&+ \sum_{i=1}^n \oint_{\partial D^i_\varepsilon} \langle n_i, \tilde v_i \rangle \sum_{j=1}^n \langle \nabla_{\! p_j} \rho_\varepsilon^\pm, v_j \rangle \dd s \label{eq:l3}\\
	&+ \sum_{j=1}^n \oint_{\partial D^j_\varepsilon} \langle n_j, \tilde v_j \rangle \sum_{i=1}^n \langle \pm 2\pi \nabla_{\! p_i} G_i, v_i \rangle \rho_\varepsilon^\pm \, \dd s \label{eq:l4}\\
	&+ \sum_{i=1}^n \oint_{\partial D^i_\varepsilon} \langle n_i, \tilde v_i \rangle \langle \nabla_{\! z} \rho_\varepsilon^\pm, \tilde v_i \rangle \dd s \nonumber
\end{align}

We expand each term
\begin{flalign*}
	\text{\eqref{eq:l1}}
	=& \int_{\Omega_\varepsilon} \sum_{i=1}^n \pm 2\pi \Hess_{\! z} G_i (\tilde v_i, \tilde v_i) \dd\mu_\varepsilon^\pm &&
\end{flalign*}
\begin{flalign*}
	\text{\eqref{eq:l2}}
	=& \int_{\Omega_\varepsilon} \sum_{i=1}^n \langle \mp 2\pi \nabla_{\! z} G_i, \tilde v_i \rangle \sum_{j=1}^n \langle \nabla_{\! p_j} \rho_\varepsilon^\pm, v_j \rangle \, \dd A_\sph &&\\
	=& \int_{\Omega_\varepsilon} \sum_{i=1}^n \langle \mp 2\pi \nabla_{\! z} G_i, \tilde v_i \rangle \sum_{j=1}^n \langle \mp 2\pi \nabla_{\! z} G_j, \tilde v_j \rangle \dd\mu_\varepsilon^\pm &&\\
	&- \int_{\Omega_\varepsilon} \sum_{i=1}^n \langle \mp 2\pi \nabla_{\! z} G_i, \tilde v_i \rangle \dd\mu_\varepsilon^\pm
	\int_{\Omega_\varepsilon} \sum_{j=1}^n \langle \mp 2\pi \nabla_{\! z} G_j, \tilde v_j \rangle \dd\mu_\varepsilon^\pm &&\\
	&- \int_{\Omega_\varepsilon} \sum_{i=1}^n \langle \mp 2\pi \nabla_{\! z} G_i, \tilde v_i \rangle \dd\mu_\varepsilon^\pm 
	\sum_{j=1}^n \oint_{\partial D^j_\varepsilon} \langle n_j, \tilde v_j \rangle \rho_\varepsilon^\pm \, \dd s &&
\end{flalign*}
\begin{flalign*}
	\text{\eqref{eq:l3}}
	=& \sum_{i=1}^n \oint_{\partial D^i_\varepsilon} \langle n_i, \tilde v_i \rangle \sum_{j=1}^n \langle \mp 2\pi \nabla_{\! z} G_j, \tilde v_j \rangle \rho_\varepsilon^\pm \, \dd s &&\\
	&- \sum_{i=1}^n \oint_{\partial D^i_\varepsilon} \langle n_i, \tilde v_i \rangle\rho_\varepsilon^\pm \, \dd s
	\int_{\Omega_\varepsilon} \sum_{j=1}^n \langle \mp 2\pi \nabla_{\! z} G_j, \tilde v_j \rangle \dd\mu_\varepsilon^\pm &&\\
	&- \sum_{i=1}^n \oint_{\partial D^i_\varepsilon} \langle n_i, \tilde v_i \rangle\rho_\varepsilon^\pm \, \dd s
	\sum_{j=1}^n \oint_{\partial D^j_\varepsilon} \langle n_j, \tilde v_j \rangle \rho_\varepsilon^\pm \, \dd s, &&
\end{flalign*}
and
\begin{flalign*}
	\text{\eqref{eq:l4}}
	=& \sum_{j=1}^n \int_{\partial D^j_\varepsilon} \langle n_j, \tilde v_j \rangle \sum_{i=1}^n \langle \mp 2\pi \nabla_{\! z} G_i, \tilde v_i \rangle \rho_\varepsilon^\pm \, \dd s. &&
\end{flalign*}

Putting them together gives
\begin{align}
	\Hess_{\! p} \log I^\pm_\varepsilon (v, v)
	=& \int_{\Omega_\varepsilon} \Big( \sum_{i=1}^n \langle 2\pi \nabla_{\! z} G_i, \tilde v_i \rangle \Big)^2 \dd\mu_\varepsilon^\pm \label{eq:h1}\\
&+ \int_{\Omega_\varepsilon} \sum_{i=1}^n \pm 2\pi \Hess_{\! z} G_i (\tilde v_i, \tilde v_i) \dd\mu_\varepsilon^\pm  \label{eq:h2}\\
	&- \Big(
		\int_{\Omega_\varepsilon} \sum_{i=1}^n \langle \mp 2\pi \nabla_{\! z} G_i, \tilde v_i \rangle \dd\mu_\varepsilon^\pm
		+ \sum_{i=1}^n \oint_{\partial D^i_\varepsilon} \langle n_i, \tilde v_i \rangle \rho_\varepsilon^\pm \dd s
	\Big)^2 \label{eq:h3}\\
	&+ \sum_{i=1}^n \oint_{\partial D^i_\varepsilon} \langle n_i, \tilde v_i \rangle \sum_{j=1}^n \langle \mp 2\pi \nabla_{\! z} G_j, \tilde v_j \rangle \rho_\varepsilon^\pm \, \dd s  \label{eq:h4}\\
	&+ \sum_{i=1}^n \oint_{\partial D^i_\varepsilon} \langle n_i, \tilde v_i \rangle \sum_{j=1}^n \langle \mp 2\pi \nabla_{\! z} G_j, \tilde v_j - \tilde v_i \rangle \rho_\varepsilon^\pm \, \dd s \label{eq:h5}
\end{align}

Note again the independence to the choice of the Killing fields $\tilde v_i$
(by the same argument as before) and, when $\tilde v_i = \tilde v_*$ is the
same Killing field for all $i$,
\begin{multline*}
	\Hess_{\! p} \log I^\pm_\varepsilon (v, v)
	= \int_{\Omega_\varepsilon} \langle \nabla_{\! z} u, \tilde v_* \rangle^2 \dd\mu_\varepsilon^\pm
	+ \int_{\Omega_\varepsilon} \pm \Hess_{\! z} u (\tilde v_*, \tilde v_*) \dd\mu_\varepsilon^\pm\\
	+ \sum_{i=1}^n \oint_{\partial D^i_\varepsilon} \langle n_i, \tilde v_* \rangle \langle \mp \nabla_{\! z} u, \tilde v_* \rangle \rho_\varepsilon^\pm \, \dd s = 0.
\end{multline*}

The boundary term~\eqref{eq:h5} converges to $0$ because
$D^i_\varepsilon$ only contains the singularity of $G_j$ when $j=i$.  For the
same reason, in the limit $\varepsilon \to 0$, we may rewrite the ``coupling''
boundary term \eqref{eq:h4} into the form
\begin{align*}
	&\lim_{\varepsilon \to 0} \sum_{i=1}^n \oint_{\partial D^i_\varepsilon} \langle n_i, \tilde v_i \rangle \langle \mp 2\pi \nabla_{\! z} G_i, \tilde v_i \rangle \rho_\varepsilon^\pm \, \dd s\\
	=&\lim_{\varepsilon \to 0} \sum_{i=1}^n \Big[
		- \int_{\Omega_\varepsilon} \langle \nabla_{\! z} u, \tilde v_i \rangle \langle 2\pi \nabla_{\! z} G_i, \tilde v_i \rangle \, \dd \mu_\varepsilon^\pm
		\mp \int_{\Omega_\varepsilon} 2\pi \Hess_{\! z} G_i (\tilde v_i, \tilde v_i) \, \dd \mu_\varepsilon^\pm
	\Big],
\end{align*}
where the second term cancels with~\eqref{eq:h2}.

We have seen that the second term in~\eqref{eq:h3} converges to $0$ as
$\varepsilon \to 0$.  At a critical point, its first term~\eqref{eq:h3}
vanishes.  If this is the case, we have
\begin{align*}
	&\lim_{\varepsilon \to 0} \Hess_{\! p} \log I^\pm_\varepsilon (v, v) \\
	=&\lim_{\varepsilon \to 0}
	\int_{\Omega_\varepsilon} \Big( \sum_{i=1}^n \langle 2\pi \nabla_{\! z} G_i, \tilde v_i \rangle \Big)^2 \dd\mu_\varepsilon^\pm
	- \int_{\Omega_\varepsilon} \sum_{i=1}^n \langle \nabla_{\! z} u, \tilde v_i \rangle \langle 2\pi \nabla_{\! z} G_i, \tilde v_i \rangle \, \dd \mu_\varepsilon^\pm\\
	=&\lim_{\varepsilon \to 0}
	\int_{\Omega_\varepsilon} \sum_{i=1}^n \sum_{j=1}^n \langle 2\pi \nabla_{\! z} G_j, \tilde v_j - \tilde v_i \rangle \langle 2\pi \nabla_{\! z} G_i, \tilde v_i \rangle \, \dd \mu_\varepsilon^\pm.
\end{align*}
When $i \ne j$, the convergence follows from the same argument for the
convergence of the gradience~\eqref{eq:grad}.  The possible divergence with $i
= j$ is ruled out by $\tilde v_j - \tilde v_i$.  In the following, we will
simply write
\[
	\Hess_{\! p} \log I^\pm (v, v)
	=
	\int_{\uS^2} \sum_{i=1}^n \sum_{j=1}^n \langle 2\pi \nabla_{\! z} G_j, \tilde v_j - \tilde v_i \rangle \langle 2\pi \nabla_{\! z} G_i, \tilde v_i \rangle \, \dd \mu^\pm.
\]

\section{Local minimizer}

Without loss of generality, we assume that the vertices of the cube, in the
standard coordinates $(\theta, \phi)$ of the sphere, are as follows:
\begin{align*}
	p_1 &= (\theta_0, 0),&
	p_5 &= (\pi - \theta_0, 0),\\
	p_2 &= (\theta_0, \pi/2),&
	p_6 &= (\pi - \theta_0, \pi/2),\\
	p_3 &= (\theta_0, \pi),&
	p_7 &= (\pi - \theta_0, \pi),\\
	p_4 &= (\theta_0, 3\pi/2),&
	p_8 &= (\pi - \theta_0, 3\pi/2).
\end{align*}
where $\theta_0 = \arccos(1/\sqrt{3})$.

\subsection{Decomposition into irreducible representations}

We have seen that the Hessian vanishes on Killing fields, compatible with the
fact the $J$ is invariant under global rotations.  So we study the Hessian on
the $13$-dimensional space
\[
	V = \Big(\bigoplus_{i=1}^8 T_{p_i} \uS^2\Big) / \mathrm{SO}(3)
\]
at the cubic configuration.  For this purpose, we decompose $V$ into
irreducible representations (irreps) of the symmetry group of the cube, as
follows:

\begin{itemize}
	% \item The $3$-dimensional space $T_{1g}$ generated by rotations.  One
	% 	generator is given by
	% 	\[
	% 		v_i = \frac{\partial_\phi}{\sin\theta}\Big|_{p_i} \quad 1 \le i \le 8.
	% 	\]
	\item The $3$-dimensional space $T_{1u}$ generated by slides (that moves the
		center of mass).  One generator is given by
		\[
			% v_i = \partial_\theta|_{p_i}, \quad 1 \le i \le 8.
			v_i = \frac{\partial_\phi}{\sin\theta}\Big|_{p_i} \text{ for } i \in \{1, 2, 5, 6\},\qquad
			v_i =-\frac{\partial_\phi}{\sin\theta}\Big|_{p_i} \text{ for } i \in \{3, 4, 7, 8\}.
		\]
	\item The $2$-dimensional space $E_g$ generated by orbit of
		\[
			v_i = \frac{\partial_\phi}{\sin\theta}\Big|_{p_i} \text{ for } i \in \{1, 3, 5, 7\},\qquad
			v_i =-\frac{\partial_\phi}{\sin\theta}\Big|_{p_i} \text{ for } i \in \{2, 4, 6, 8\}.
		\]
		Every face remain rectangular under these deformations.
	\item The $2$-dimensional space $E_u$ generated by orbit of
		\[
			v_i = \frac{\partial_\phi}{\sin\theta}\Big|_{p_i} \text{ for } i \in \{1, 2, 3, 4\},\qquad
			v_i =-\frac{\partial_\phi}{\sin\theta}\Big|_{p_i} \text{ for } i \in \{5, 6, 7, 8\}.
		\]
		The face centers of the cube remain symmetry centers under these deformations.
	\item The $3$-dimensional space $T_{2g}$ generated by the orbit of
		\[
			% v_i = \partial_\theta|_{p_i}, & \text{ for } i \in \{1, 3, 6, 8\},\qquad
			% v_i =-\partial_\theta|_{p_i}, & \text{ for } i \in \{2, 4, 5, 7\}.
			v_i = \frac{\partial_\phi}{\sin\theta}\Big|_{p_i} \text{ for } i \in \{1, 2, 7, 8\},\qquad
			v_i =-\frac{\partial_\phi}{\sin\theta}\Big|_{p_i} \text{ for } i \in \{3, 4, 5, 6\}.
		\]
	\item The $3$-dimensional space $T_{2u}$ generated by orbit of
		\[
			v_i = \frac{\partial_\phi}{\sin\theta}\Big|_{p_i} \text{ for } i \in \{1, 3, 6, 8\},\qquad
			v_i =-\frac{\partial_\phi}{\sin\theta}\Big|_{p_i} \text{ for } i \in \{2, 4, 5, 7\}.
		\]
\end{itemize}

The Hessian is invariant under the action of $G$.  By Schur's lemma, it must be
a multiple of identity on each of the irreps.  So we only need to verify that
the Hessian is non-negative on one generator for each of the irreps.

\subsection{The definiteness on \texorpdfstring{$T_{1u}$}{T1u}, \texorpdfstring{$E_u$}{Eu} and \texorpdfstring{$E_g$}{Eg}}

We have chosen the generator vectors, intentionally, with the following
property.  The eight vertices are partitioned into two subsets $A$ and $B$,
with four vertices each, such that vertices in $A$ extend to the Killing field
$\partial_\phi$, while vertices in $B$ extend to $-\partial_\phi$.  In this
case, we have
\begin{align*}
	&\Hess_{\! p} \log I^\pm (v, v)
	=
	\int_{\uS^2} \sum_{i=1}^n \sum_{j=1}^n \langle 2\pi \nabla_{\! z} G_j, \tilde v_j - \tilde v_i \rangle \langle 2\pi \nabla_{\! z} G_i, \tilde v_i \rangle \, \dd \mu^\pm\\
	=& -4
	\int_{\uS^2} \sum_{i \in A} \sum_{j \in B} (2\pi \partial_\phi G_j)(2\pi \partial_\phi G_i) \, \dd \mu^\pm
	= -16\pi^2 \int_{\uS^2} \partial_\phi G_A \partial_\phi G_B \, \dd \mu^\pm,
\end{align*}
where
\[
	G_A = \sum_{i \in A} G_i, \qquad G_B = \sum_{i \in B} G_i.
\]

To study the definiteness of the Hessian, we need to first study the sign of
\[
	H^\pm_{ij} = \int_{\uS^2} \partial_\phi G_i \partial_\phi G_j \,
	\dd \mu^\pm.
\]
When $p_i$ and $p_j$ are antipodal (e.g. $(i,j) = (1,7)$) or on a horizontal
face diagonal (e.g. $(i,j) = (1,3)$), one easily sees that the integrand
$\partial_\phi G_i \partial_\phi G_j \le 0$ everywhere on the sphere, so
$H^\pm_{ij}$ is negative.  When $p_i$ and $p_j$ are vertical neighbors (e.g.
$(i,j) = (1,5)$), one easily sees that the integrand $\partial_\phi G_i
\partial_\phi G_j \ge 0$ everywhere on the sphere, so $H^\pm_{ij}$ is positive.
Otherwise, the following lemma would be useful:

\begin{lemma}
	Let $f(\theta, \phi)$ be a function defined on $\uS^2$ that is odd in $\phi$,
	and $\rho$ be a positive function on $\uS^2$ with
	\[
		\rho(\theta, \phi) = \rho(\theta, \phi+\tfrac{\pi}{2}) = \rho(\pi - \theta, \phi).
	\]
	Assume that $f(\theta, \phi)$ is negative and strictly monotonically
	increasing in $\phi$ on the interval $(0, \pi)$.  Then the integrals
	\begin{align*}
		K_1 =& \int_{\uS^2} f(\theta, \phi) f(\theta, \phi - \tfrac{\pi}{2}) \rho \, dA_\sph < 0,\\
		K_2 =& \int_{\uS^2} f(\theta, \phi) f(\pi - \theta, \phi - \tfrac{\pi}{2}) \rho \, dA_\sph < 0.
	\end{align*}
	Moreover, if $f(\theta, \phi) - f(\pi - \theta, \phi)$ is negative and
	monotonically increasing in $\phi$ on the interval $(0, \pi)$ for $0 < \theta
	< \tfrac{\pi}{2}$, then $K_1 < K_2 < 0$.
\end{lemma}

\begin{proof}
	We decompose the sphere into four quadrants.  Then
	\begin{align*}
		K_1 = \int_0^\pi \sin\theta \, d\theta \int_{0}^{\tfrac{\pi}{2}} & \rho \, d\phi
		\big(f(\theta, \phi) f(\theta, \phi-\tfrac{\pi}{2}) + f(\theta, \phi+\tfrac{\pi}{2}) f(\theta, \phi)\\
		&+ f(\theta, \phi+\pi) f(\theta, \phi+\tfrac{\pi}{2}) + f(\theta, \phi-\tfrac{\pi}{2}) f(\theta, \phi+\pi) \big)\\
		= \int_0^\pi \sin\theta \, d\theta \int_{0}^{\tfrac{\pi}{2}} \rho \, d\phi
		& \big( f(\theta, \phi) + f(\theta, \phi + \pi) \big) \big( f(\theta, \phi - \tfrac{\pi}{2}) + f(\theta, \phi + \tfrac{\pi}{2}) \big)\\
		= \int_0^\pi \sin\theta \, d\theta \int_{0}^{\tfrac{\pi}{2}} \rho \, d\phi
		& \big( f(\theta, \phi) - f(\theta, \pi - \phi) \big) \big(-f(\theta, \tfrac{\pi}{2} - \phi) + f(\theta, \tfrac{\pi}{2} + \phi) \big)
	\end{align*}
	For $0 < \phi < \tfrac{\pi}{2}$, by the assumed monotonicity, we have $f(\phi) < f(\pi
	- \phi) < 0$ and $g(\tfrac{\pi}{2}-\phi) < g(\tfrac{\pi}{2} + \phi) < 0$.  So $K_1$ is
	negative.  On the other hand,
	\begin{align*}
		K_2 = \int_0^{\tfrac{\pi}{2}} & \sin\theta \, d\theta \int_{0}^{\tfrac{\pi}{2}} \rho \, d\phi \big(
			f(\theta, \phi) f(\pi - \theta, \phi-\tfrac{\pi}{2}) + f(\pi - \theta, \phi) f(\theta, \phi-\tfrac{\pi}{2})\\
			&+f(\theta, \phi+\tfrac{\pi}{2}) f(\pi - \theta, \phi) + f(\pi - \theta, \phi+\tfrac{\pi}{2}) f(\theta, \phi)\\
			&+f(\theta, \phi+\pi) f(\pi - \theta, \phi+\tfrac{\pi}{2}) + f(\pi - \theta, \phi+\pi) f(\theta, \phi+\tfrac{\pi}{2})\\
			&+f(\theta, \phi-\tfrac{\pi}{2}) f(\pi - \theta, \phi+\pi) + f(\pi - \theta, \phi-\tfrac{\pi}{2}) f(\theta, \phi+\pi)
		\big)\\
		= \int_0^{\tfrac{\pi}{2}} & \sin\theta \, d\theta \int_{0}^{\tfrac{\pi}{2}} \rho \, d\phi \\ & \Big[
			\big(f(\pi - \theta, \phi) + f(\pi - \theta, \phi + \pi) \big) \big( f(\theta, \phi - \tfrac{\pi}{2}) + f(\theta, \phi + \tfrac{\pi}{2}) \big) \\
			&+\big(f(\theta, \phi) + f(\theta, \phi + \pi) \big) \big( f(\pi - \theta, \phi - \tfrac{\pi}{2}) + f(\pi - \theta, \phi + \tfrac{\pi}{2}) \big) 
		\Big] \\
		= \int_0^{\tfrac{\pi}{2}} & \sin\theta \, d\theta \int_{0}^{\tfrac{\pi}{2}}\rho \, d\phi \\ & \Big[
			\big(f(\pi - \theta, \phi) - f(\pi - \theta, \pi - \phi) \big) \big( -f(\theta, \tfrac{\pi}{2} - \phi) + f(\theta, \phi + \tfrac{\pi}{2}) \big) \\
			&+\big(f(\theta, \phi) - f(\theta, \pi - \phi) \big) \big( -f(\pi - \theta, \tfrac{\pi}{2} - \phi) + f(\pi - \theta, \phi + \tfrac{\pi}{2}) \big) 
		\Big]
	\end{align*}
	and the integrand is negative by the same argument as above.

	Moreover,
	\begin{align*}
		K_1 - K_2 & = \int_0^{\tfrac{\pi}{2}} \sin\theta \, d\theta \int_{0}^{\tfrac{\pi}{2}} \rho \, d\phi \\ &
		\big(
			f(\theta,\phi) - f(\pi - \theta, \phi) - f(\theta, \pi - \phi) + f(\pi - \theta, \pi - \phi)
		\big) \\ &
		\big(
			-f(\theta,\tfrac{\pi}{2}-\phi) + f(\pi - \theta, \tfrac{\pi}{2} - \phi) + f(\theta, \tfrac{\pi}{2} + \phi) - f(\pi - \theta, \tfrac{\pi}{2} + \phi)
		\big).
	\end{align*}
	Again the integrand is negative if $f(\theta, \phi) - f(\pi - \theta, \phi)$
	is negative and monotonically increasing in $\phi$ on the interval $(0,
	\pi)$.  So $K_1 < K_2$.
\end{proof}

Using this lemma, one easily proves that if $p_i$ and $p_j$ are horizontal
neighbors (e.g. $(i,j) = (1,2)$) or on a vertical face diagonal (e.g. $(i,j) =
(1,6)$), we have $H^\pm_{ij} < 0$.  Consequently, by
\[
	\Hess_{\! p} \log I^\pm (v, v) = -16 \pi^2 \sum_{i \in A} \sum_{j \in B} H^\pm_{ij}
\]
we can already conclude that the Hessian is positive definite on
the irreps $T_{1u}$, $E_g$, but negative definite on $E_u$.

\subsection{The definiteness on \texorpdfstring{$T_{2u}$}{T2u}}

For $T_{2u}$ we place the cube so that a pair of antipodal vertices are at the
north and the south pole.  So
\begin{align*}
	p_1 &= (0, 0),&
	p_5 &= (\pi, 0),\\
	p_2 &= (\theta_0, 0),&
	p_6 &= (\pi - \theta_0, \pi/3),\\
	p_3 &= (\theta_0, 2\pi/3),&
	p_7 &= (\pi - \theta_0, \pi),\\
	p_4 &= (\theta_0, 4\pi/3),&
	p_8 &= (\pi - \theta_0, 5\pi/3).
\end{align*}
where $\theta_0 = \arccos(-1/3)$.  Then a generator for $T_{2u}$ is given by
\[
	v_i = \frac{\partial_\phi}{\sin\theta}\Big|_{p_i} \text{ for } i \in \{1, 2, 3, 4\},\qquad
	v_i =-\frac{\partial_\phi}{\sin\theta}\Big|_{p_i} \text{ for } i \in \{5, 6, 7, 8\}.
\]
Note that the vertices $p_1$ and $p_5$ (at the poles) are actually fixed under
this generator.  The Hessian becomes
\[
	\Hess_{\! p} \log I^\pm (v, v)
	= -16\pi^2 \int_{\uS^2} \partial_\phi G_A \partial_\phi G_B \, \dd \mu^\pm,
\]
where $A=\{2,3,4\}$ and $B=\{6,7,8\}$.  One easily notices that the integrand
$\partial_\phi G_A \partial_\phi G_B \le 0$ everywhere on the sphere.  So the
Hessian is positive definite on $T_{2u}$.

\subsection{The definiteness on \texorpdfstring{$T_{2g}$}{T2g}}

By Schur's lemma, $\Hess|_{T_{2g}}=\lambda\cdot\mathrm{Id}$ for some scalar
$\lambda$.  Since $J_8$ is invariant under $\SO(3)$ acting simultaneously on all
eight vertices, the eigenvalue $\lambda$ is the same for any cubic configuration.
We determine the sign of $\lambda$ by evaluating the Hessian on a convenient
$T_{2g}$-generator.

\paragraph{A convenient cube orientation.}
Set $a:=\sqrt{2/3}$ and $b:=1/\sqrt{3}$, and place the cube with vertex sets
\[
	A':=\bigl\{(\pm b,\,0,\,\pm a)\bigr\},
	\qquad
	B':=\bigl\{(\pm b,\,\pm a,\,0)\bigr\}.
\]
These eight points are vertices of a regular cube on $\uS^2$ with edge length
$2/\sqrt{3}$: for instance, the face $\{X=b\}$ has vertices
$(b,0,a),(b,a,0),(b,0,-a),(b,-a,0)$ forming a square of side $a\sqrt{2}=2/\sqrt{3}$.

Let $G_{A'}:=\sum_{p_i\in A'}G_{p_i}$, $G_{B'}:=\sum_{p_i\in B'}G_{p_i}$, and
let $\mu^\pm$ denote the probability measures for this configuration.

\paragraph{\texorpdfstring{$T_{2g}$}{T2g} generator.}
Let $\partial_\psi:=Y\partial_Z-Z\partial_Y$ be the Killing field for rotation
about the $X$-axis, and define
\[
	w_{p_i}:=+\partial_\psi\big|_{p_i}\;\text{ for }p_i\in A',
	\qquad
	w_{p_i}:=-\partial_\psi\big|_{p_i}\;\text{ for }p_i\in B'.
\]
Every vertex satisfies $\sqrt{Y_i^2+Z_i^2}=a$, so each $w_{p_i}$ has magnitude $a$.
The vector $w$ is a $T_{2g}$ generator: the cube symmetry
$R_X\colon(X,Y,Z)\mapsto(X,-Z,Y)$ (rotation by $90^\circ$ about $X$) maps
$A'\to B'$ and preserves $\partial_\psi$, so $w\mapsto -w$ under $R_X$,
in accordance with the $T_{2g}$ character $(-1)$ under a $C_4$ rotation.
Both the partition $\{A',B'\}$ and the Killing field $\partial_\psi$ are even
under the point inversion $(X,Y,Z)\mapsto(-X,-Y,-Z)$, confirming $w\in T_{2g}$
(a $g$-type irrep).

The general Hessian formula (with $\partial_\psi$ in place of $\partial_\phi$) gives
\[
	\Hess_p\log J_8(w,w)
	= -16\pi^2\int_{\uS^2}\partial_\psi G_{A'}\,\partial_\psi G_{B'}\,\dd(\mu^++\mu^-).
\]
It remains to show this integral is strictly negative.

\paragraph{Computing $\partial_\psi G_{A'}$ and $\partial_\psi G_{B'}$.}
The $A'$-vertices lie in the $XZ$-plane ($Y_i=0$), so the inner products
$P\cdot p_i = \pm bX\pm aZ$ depend only on $X$ and $Z$.  Hence
\[
	G_{A'}(X,Y,Z)
	= -\frac{1}{4\pi}\log P(X,Z)+C_{A'},
	\qquad
	P(s,t):=\bigl((1-at)^2-b^2s^2\bigr)\bigl((1+at)^2-b^2s^2\bigr).
\]
Since $P(X,Z)$ does not depend on $Y$, and $\partial_\psi Z = Y$:
\[
	\partial_\psi G_{A'}
	= -\frac{Y\,\partial_t P(X,Z)}{4\pi\,P(X,Z)}.
\]
A direct computation yields $\partial_t P(s,t)=-4a^2t\,(1+b^2s^2-a^2t^2)$, and therefore
\[
	\partial_\psi G_{A'}
	= \frac{a^2\,YZ\,(1+b^2X^2-a^2Z^2)}{\pi\,P(X,Z)}.
\]
Analogously, the $B'$-vertices lie in the $XY$-plane ($Z_i=0$), giving
$G_{B'}=-\tfrac{1}{4\pi}\log P(X,Y)+C_{B'}$, and since $\partial_\psi[P(X,Y)]=-Z\,\partial_t P(X,Y)$:
\[
	\partial_\psi G_{B'}
	= \frac{-a^2\,YZ\,(1+b^2X^2-a^2Y^2)}{\pi\,P(X,Y)}.
\]

\paragraph{The integrand is non-positive.}
Multiplying,
\[
	\partial_\psi G_{A'}\cdot\partial_\psi G_{B'}
	= \frac{-a^4\,Y^2Z^2\,(1+b^2X^2-a^2Z^2)(1+b^2X^2-a^2Y^2)}{\pi^2\,P(X,Z)\,P(X,Y)}.
\]
On $\uS^2$ (using $a^2=2/3$, $b^2=1/3$, and $X^2+Y^2+Z^2=1$):
\begin{align*}
	1+b^2X^2-a^2Z^2
	&= \tfrac{4}{3}X^2+Y^2+\tfrac{1}{3}Z^2 > 0,\\
	1+b^2X^2-a^2Y^2
	&= \tfrac{4}{3}X^2+\tfrac{1}{3}Y^2+Z^2 > 0.
\end{align*}
The denominators $P(X,Z)$ and $P(X,Y)$ are also positive away from the cube
vertices: on $\uS^2$ one checks that
\[
	(1-aZ)^2-b^2X^2 = (Z-a)^2+\tfrac13 Y^2\ge 0,
\]
vanishing only when $Z=a$ and $Y=0$, i.e.\ at the $A'$-vertices; similarly for
the other factor and for $P(X,Y)$ at the $B'$-vertices.
Since $Y^2Z^2\ge0$ and both factors in the numerator are strictly positive, the
entire product $\partial_\psi G_{A'}\cdot\partial_\psi G_{B'}$ is non-positive,
and vanishes only when $Y=0$ or $Z=0$---a set of zero spherical measure.

Since $\mu^++\mu^-$ is absolutely continuous with respect to the spherical area
measure,
\[
	\int_{\uS^2}\partial_\psi G_{A'}\,\partial_\psi G_{B'}\,\dd(\mu^++\mu^-)
	< 0.
\]
Therefore $\Hess_p\log J_8(w,w)>0$.  Since the $T_{2g}$ eigenvalue $\lambda$ is
the same for any cube orientation, the Hessian is positive definite on $T_{2g}$.

\bibliography{References}
\bibliographystyle{plain}

\end{document}
